\documentclass[a4paper,11pt,oneside]{memoir}

% Castellano
\usepackage[spanish]{babel}
\selectlanguage{spanish}
\usepackage[utf8]{inputenc}
\usepackage{placeins}
\usepackage{hologo}
\usepackage{colortbl}


\RequirePackage{booktabs}
\RequirePackage[table]{xcolor}
\RequirePackage{xtab}
\RequirePackage{multirow}

% Links
\usepackage[colorlinks]{hyperref}
\hypersetup{
	allcolors = {red}
}

% Ecuaciones
\usepackage{amsmath}

% Rutas de fichero / paquete
\newcommand{\ruta}[1]{{\sffamily #1}}

% Párrafos
\nonzeroparskip


% Imagenes
\usepackage{graphicx}
\newcommand{\imagen}[2]{
	\begin{figure}[!h]
		\centering
		\includegraphics[width=0.9\textwidth]{#1}
		\caption{#2}\label{fig:#1}
	\end{figure}
	\FloatBarrier
}

\newcommand{\imagenflotante}[2]{
	\begin{figure}%[!h]
		\centering
		\includegraphics[width=0.9\textwidth]{#1}
		\caption{#2}\label{fig:#1}
	\end{figure}
}



% El comando \figura nos permite insertar figuras comodamente, y utilizando
% siempre el mismo formato. Los parametros son:
% 1 -> Porcentaje del ancho de página que ocupará la figura (de 0 a 1)
% 2 --> Fichero de la imagen
% 3 --> Texto a pie de imagen
% 4 --> Etiqueta (label) para referencias
% 5 --> Opciones que queramos pasarle al \includegraphics
% 6 --> Opciones de posicionamiento a pasarle a \begin{figure}
\newcommand{\figuraConPosicion}[6]{%
  \setlength{\anchoFloat}{#1\textwidth}%
  \addtolength{\anchoFloat}{-4\fboxsep}%
  \setlength{\anchoFigura}{\anchoFloat}%
  \begin{figure}[#6]
    \begin{center}%
      \Ovalbox{%
        \begin{minipage}{\anchoFloat}%
          \begin{center}%
            \includegraphics[width=\anchoFigura,#5]{#2}%
            \caption{#3}%
            \label{#4}%
          \end{center}%
        \end{minipage}
      }%
    \end{center}%
  \end{figure}%
}

%
% Comando para incluir imágenes en formato apaisado (sin marco).
\newcommand{\figuraApaisadaSinMarco}[5]{%
  \begin{figure}%
    \begin{center}%
    \includegraphics[angle=90,height=#1\textheight,#5]{#2}%
    \caption{#3}%
    \label{#4}%
    \end{center}%
  \end{figure}%
}
% Para las tablas
\newcommand{\otoprule}{\midrule [\heavyrulewidth]}
%
% Nuevo comando para tablas pequeñas (menos de una página).
\newcommand{\tablaSmall}[5]{%
 \begin{table}
  \begin{center}
   \rowcolors {2}{gray!35}{}
   \begin{tabular}{#2}
    \toprule
    #4
    \otoprule
    #5
    \bottomrule
   \end{tabular}
   \caption{#1}
   \label{tabla:#3}
  \end{center}
 \end{table}
}

%
% Nuevo comando para tablas pequeñas (menos de una página).
\newcommand{\tablaSmallSinColores}[5]{%
 \begin{table}[H]
  \begin{center}
   \begin{tabular}{#2}
    \toprule
    #4
    \otoprule
    #5
    \bottomrule
   \end{tabular}
   \caption{#1}
   \label{tabla:#3}
  \end{center}
 \end{table}
}

\newcommand{\tablaApaisadaSmall}[5]{%
\begin{landscape}
  \begin{table}
   \begin{center}
    \rowcolors {2}{gray!35}{}
    \begin{tabular}{#2}
     \toprule
     #4
     \otoprule
     #5
     \bottomrule
    \end{tabular}
    \caption{#1}
    \label{tabla:#3}
   \end{center}
  \end{table}
\end{landscape}
}

%
% Nuevo comando para tablas grandes con cabecera y filas alternas coloreadas en gris.
\newcommand{\tabla}[6]{%
  \begin{center}
    \tablefirsthead{
      \toprule
      #5
      \otoprule
    }
    \tablehead{
      \multicolumn{#3}{l}{\small\sl continúa desde la página anterior}\\
      \toprule
      #5
      \otoprule
    }
    \tabletail{
      \hline
      \multicolumn{#3}{r}{\small\sl continúa en la página siguiente}\\
    }
    \tablelasttail{
      \hline
    }
    \bottomcaption{#1}
    \rowcolors {2}{gray!35}{}
    \begin{xtabular}{#2}
      #6
      \bottomrule
    \end{xtabular}
    \label{tabla:#4}
  \end{center}
}

%
% Nuevo comando para tablas grandes con cabecera.
\newcommand{\tablaSinColores}[6]{%
  \begin{center}
    \tablefirsthead{
      \toprule
      #5
      \otoprule
    }
    \tablehead{
      \multicolumn{#3}{l}{\small\sl continúa desde la página anterior}\\
      \toprule
      #5
      \otoprule
    }
    \tabletail{
      \hline
      \multicolumn{#3}{r}{\small\sl continúa en la página siguiente}\\
    }
    \tablelasttail{
      \hline
    }
    \bottomcaption{#1}
    \begin{xtabular}{#2}
      #6
      \bottomrule
    \end{xtabular}
    \label{tabla:#4}
  \end{center}
}

%
% Nuevo comando para tablas grandes sin cabecera.
\newcommand{\tablaSinCabecera}[5]{%
  \begin{center}
    \tablefirsthead{
      \toprule
    }
    \tablehead{
      \multicolumn{#3}{l}{\small\sl continúa desde la página anterior}\\
      \hline
    }
    \tabletail{
      \hline
      \multicolumn{#3}{r}{\small\sl continúa en la página siguiente}\\
    }
    \tablelasttail{
      \hline
    }
    \bottomcaption{#1}
  \begin{xtabular}{#2}
    #5
   \bottomrule
  \end{xtabular}
  \label{tabla:#4}
  \end{center}
}



\definecolor{cgoLight}{HTML}{EEEEEE}
\definecolor{cgoExtralight}{HTML}{FFFFFF}

%
% Nuevo comando para tablas grandes sin cabecera.
\newcommand{\tablaSinCabeceraConBandas}[5]{%
  \begin{center}
    \tablefirsthead{
      \toprule
    }
    \tablehead{
      \multicolumn{#3}{l}{\small\sl continúa desde la página anterior}\\
      \hline
    }
    \tabletail{
      \hline
      \multicolumn{#3}{r}{\small\sl continúa en la página siguiente}\\
    }
    \tablelasttail{
      \hline
    }
    \bottomcaption{#1}
    \rowcolors[]{1}{cgoExtralight}{cgoLight}

  \begin{xtabular}{#2}
    #5
   \bottomrule
  \end{xtabular}
  \label{tabla:#4}
  \end{center}
}

















\graphicspath{ {./img/} }

% Capítulos
\chapterstyle{bianchi}
\newcommand{\capitulo}[2]{
	\setcounter{chapter}{#1}
	\setcounter{section}{0}
	\chapter*{#2}
	\addcontentsline{toc}{chapter}{#2}
	\markboth{#2}{#2}
}

% Apéndices
\renewcommand{\appendixname}{Apéndice}
\renewcommand*\cftappendixname{\appendixname}

\newcommand{\apendice}[1]{
	%\renewcommand{\thechapter}{A}
	\chapter{#1}
}

\renewcommand*\cftappendixname{\appendixname\ }

% Formato de portada
\makeatletter
\usepackage{xcolor}
\newcommand{\tutor}[1]{\def\@tutor{#1}}
\newcommand{\tutorAux}[1]{\def\@tutorAux{#1}}
\newcommand{\course}[1]{\def\@course{#1}}
\definecolor{cpardoBox}{HTML}{E6E6FF}
\def\maketitle{
  \null
  \thispagestyle{empty}
  % Cabecera ----------------
\noindent\includegraphics[width=\textwidth]{cabecera}\vspace{1cm}%
  \vfill
  % Título proyecto y escudo informática ----------------
  \colorbox{cpardoBox}{%
    \begin{minipage}{.8\textwidth}
      \vspace{.5cm}\Large
      \begin{center}
      \textbf{TFG del Grado en Ingeniería Informática}\vspace{.6cm}\\
      \textbf{\LARGE\@title{}}
      \end{center}
      \vspace{.2cm}
    \end{minipage}

  }%
  \hfill\begin{minipage}{.20\textwidth}
    \includegraphics[width=\textwidth]{escudoInfor}
  \end{minipage}
  \vfill
  % Datos de alumno, curso y tutores ------------------
  \begin{center}%
  {%
    \noindent\LARGE
    Presentado por \@author{}\\ 
    en Universidad de Burgos --- \@date{}\\
    Tutores: \@tutor{}\\
    \@tutorAux{}\\
  }%
  \end{center}%
  \null
  \cleardoublepage
  }
\makeatother

\newcommand{\nombre}{Nuño Basurto Hornillos} %%% cambio de comando

% Datos de portada
\title{Predicción de apuestas deportivas}
\author{Nuño Basurto Hornillos}
\tutor{Álvar Arnaiz González}
\tutorAux{Cristobal José Carmona del Jesús}
\date{23/01/2017}

\begin{document}

\maketitle



\null\cleardoublepage


%%%%%%%%%%%%%%%%%%%%%%%%%%%%%%%%%%%%%%%%%%%%%%%%%%%%%%%%%%%%%%%%%%%%%%%%%%%%%%%%%%%%%%%%
\pagestyle{empty}


\noindent\includegraphics[width=\textwidth]{cabecera}\vspace{1cm}

\noindent D. Álvar Arnaiz González y D. Cristóbal José Carmona del Jesús profesores del departamento de Ingeniería Civil, Área de Lenguajes y Sistemas Informáticos.

\noindent Expone:

\noindent Que el alumno D. \nombre, con DNI 71295798-F, ha realizado el Trabajo final de Grado en Ingeniería Informática titulado Predicción de apuestas deportivas de TFG. 

\noindent Y que dicho trabajo ha sido realizado por el alumno bajo la dirección del que suscribe, en virtud de lo cual se autoriza su presentación y defensa.

\begin{center} %\large
En Burgos, {\large \today}
\end{center}

\vfill\vfill\vfill

% Author and supervisor
\begin{minipage}{0.45\textwidth}
\begin{flushleft} %\large
Vº. Bº. del Tutor:\\[2cm]
D. Álvar Arnaiz González
\end{flushleft}
\end{minipage}
\hfill
\begin{minipage}{0.45\textwidth}
\begin{flushleft} %\large
Vº. Bº. del co-tutor:\\[2cm]
D. Cristóbal José Camona del Jesús
\end{flushleft}
\end{minipage}
\hfill

\vfill

% para casos con solo un tutor comentar lo anterior
% y descomentar lo siguiente
%Vº. Bº. del Tutor:\\[2cm]
%D. nombre tutor


\null\cleardoublepage
\null\cleardoublepage




\frontmatter

% Abstract en castellano
\renewcommand*\abstractname{Resumen}
\begin{abstract}
La gran cantidad de información disponible que nos es accesible, nos permite realizar análisis más exhaustivos con el fin de lograr valores concretos.

Este trabajo de fin de grado es una herramienta capaz de generar predicciones en partidos de fútbol, con el objetivo de proporcionar al usuario la información necesaria, para que este pueda orientarse a la hora de llevar a cabo apuestas deportivas.
\end{abstract}

\renewcommand*\abstractname{Descriptores}
\begin{abstract}
Scraping, bases de datos, Drupal, PHP, apuestas, predicción
\end{abstract}

\clearpage

% Abstract en inglés
\renewcommand*\abstractname{Abstract}
\begin{abstract}
The large amount of available information that is accessible for us, allows us to carry out more exhaustive analyzes in order to achieve concrete values.

This project is a tool capable of generating predictions in football matches, with the aim of providing to the user the necessary information, so that he can orientate himself when carrying out sports betting.
\end{abstract}

\renewcommand*\abstractname{Keywords}
\begin{abstract}
Scrapping, Bases de datos, Drupal, PHP, Backpropagation
\end{abstract}

\clearpage

% Indices
\tableofcontents

\clearpage

\listoffigures

\clearpage

%\listoftables

%\clearpage

\mainmatter
\capitulo{1}{Introducción}

El Aprendizaje automático (Machine Learning) va siendo cada vez más habitual en nuestro día a día, necesitamos que las máquinas sean capaces de aprender de su propia experiencia en vez de tener que pre-programarlas por nuestra cuenta.

A su vez vemos como las casas de apuestas son cada vez más numerosas y esto se debe a la gran popularidad que poseen. En concreto vamos a tratar con las apuestas deportivas, que nos ofrecen una infinidad de posibilidades en las que apostar. Nos vamos a centrar en los resultados finales  de los partidos de primera división.

Para el desarrollo de estas técnicas es necesario tener una cantidad de datos suficientemente grande con la que trabajar, estos datos los obtendremos mediante técnicas de web scraping.

La idea de este proyecto es unir una técnica de aprendizaje automático y las apuestas deportivas, así comprobar si realmente las redes neuronales utilizando la experiencia de partidos anteriores son capaces de realizar pronósticos acertados.





\capitulo{2}{Objetivos del proyecto}

Actualmente podemos encontrar una gran cantidad de información sobre cualquier cosa que podamos imaginar, esta información muchas veces aparece reflejada en estadísticas. El fútbol no es una excepción, en los análisis de los partidos se pueden detectar que algunos equipos tienden a perder cuando, por ejemplo tienen una menor posesión. Es este escenario donde se intenta plantear este proyecto dado que pudiendo observar las tendencias de los equipos en sus partidos y sus rachas, el algoritmo puede entender las tendencias de cada equipo.

\section{Cálculo del resultado final}

La red neuronal debe ser capaz de pronosticar el resultado de un partido de fútbol basándose en la experiencia adquirida de los pasados partidos y de las rachas calculadas. Dado que se trata de técnica de Machine Learning con aprendizaje supervisado es necesaria una clase, que será 1-X-2.

\section{Interfaz intuitiva para el usuario}

Esta información la hacemos llegar al usuario a través de un entorno sencillo donde el usuario no solo dispone de la información del resultado, sino que también le proporcionamos información de las diferentes cuotas de algunas casas de apuestas, para que de esta manera, pueda elegir en cual apostar. La manera de mostrarle os datos al usuario es sencilla, <<1>> si se trata de la victoria del equipo local, <<X>> si el partido va a acabar empate y <<2>> si la victoria va a caer del lado visitante, se ha decidido mostrar de esta manera ya que en España estamos acostumbrados a esta simbología gracias a la Quiniela.

\section{Automatización del algoritmo}
La intervención humana deber ser mínima y por este motivo es necesario automatizar el proyecto. La necesidad de automatización viene de que no puede haber una persona pendiente de cuándo debe ejecutar cada uno de los algoritmos.
Es por ello que será necesaria la creación de un servicio de Linux que sepa cuando debe ejecutar cada uno de los algoritmos.

\capitulo{3}{Conceptos teóricos}

En aquellos proyectos que necesiten para su comprensión y desarrollo de unos conceptos teóricos de una determinada materia o de un determinado dominio de conocimiento, debe existir un apartado que sintetice dichos conceptos.


\section{Scraping}
Técnica utilizada para simular  la navegación de un humano en Internet con el fin de extraer información del sitio web. Los datos pueden ser almacenados y analizados en una base de datos central o en otros lugares de almacenamiento.  Son diversas las técnicas que se pueden utilizar, en concreto se ha empleado\cite{_web_2016}

\section{Machine Learning}
Aquel proceso que le da a las computadoras la habilidad de aprender sin ser explícitamente programadas. Hay una segunda definición mucho más clara: \quote{Se dice que un programa de computación aprende de la experiencia E con respecto a una tarea T y alguna medida de rendimiento P, si el rendimiento en T, medido por P, mejora con la experiencia E}. 
El Machine Learning se divide en dos áreas principales, el aprendizaje supervisado y el no supervisado, mientras que el primero requiere una etiqueta a predecir, en el segundo no se sabe.\cite{_machine_learning_2014}

\section{Minería de datos} Es el proceso por el cual es posible detectar información procesable de conjuntos de datos, para después transformarla en una estructura comprensible y poder así utilizarla más tarde.
Tiende a confundirse con el Machine Learning, pero su principal diferencia es que mientras este último se usa para reproducir patrones conocidos y hacer predicciones basadas en patrones, la minería de datos descubre patrones desconocidos.
\section{Neuronas artificiales}
Similares a la idea de neuronas biológicas, forman parte de redes neuronales artificiales. Reciben una serie de entradas y dan una salida, la cual se ve condicionada por tres funciones: función de propagación, función de activación y función de transferencia.

\section{Redes neuronales artificiales}
Imitan el funcionamiento de las redes neuronales biológicas, donde un conjunto de neuronas artificiales trabajan unidas, a fin de resolver problemas relacionados con el reconocimiento de formas o con la predicción. Se parte de un conjunto de datos de entrada significativo para conseguir que la red aprenda las propiedades deseadas.

\section{Aprendizaje Supervisado}
Se dispone de un conjunto de ejemplos de el cual se conoce la respuesta, por lo que el objetivo es marcar una regla o correspondencia de manera que sea posible aproximar la respuesta para todos los objetos que se presenten. La salida de la función puede ser un valor numérico o una clase\cite{contributors_phpmyadmin_????}.
El uso más extendido del aprendizaje supervisado consiste en hacer predicciones a futuro basadas en comportamientos o características ya vistas en los datos que se contienen\cite{_machine_learning_2014}.

\section{Backpropagation}
También denominada propagación hacia atrás de errores, es un tipo de algoritmo con aprendizaje supervisado utilizado para el entrenamiento de las redes neuronales artificiales. Una vez aplicado un patrón, este se propaga a través del resto de capas hasta generar una salida, la cual se compara con la salida establecida como objetivo ya que se trata de aprendizaje supervisado. Se obtiene un error que se propaga hacia atrás cambiando los pesos, acercando así el algoritmo a un mejor resultado.
A medida que se entrena la red neuronal, las capas intermedias aprenden a organizarse para así reconocer diferentes características del espacio de entrada.\cite{mazur_step_2015}

\section{Preprocesamiento}
El ruido y las anomalías de las bases de datos reales suelen conducir a la extracción de patrones y reglas poco útiles. La finalidad del preprocesamiento es la creación de información homogénea.

\section{Demonios}
También llamados servicios, son un tipo especial de procesos que se ejecutan en segundo plano.
El sistema los inicia en el arranque, ejecutándose en este caso una tarea planificada y comprobando si la fecha del sistema coincide con las fechas de la jornada. Para la ejecución del demonio se emplea un gestor de servicios, llamado cron. El demonio también es posible ejecutarlo a mano ya que es posible que en un momento determinado no se encuentre el servidor iniciado y no haya sido posible la ejecución del demonio\citep{torres_servicios_2013}.
\capitulo{4}{Técnicas y herramientas}


En esta sección se explica brevemente la metodología utilizada en el desarrollo del trabajo de fin de grado. También veremos las herramientas elegidas para el llevar a cabo el mismo y su elección frente a otras.

\section{Técnicas}

\subsection{Scrum} 
Es aquel proceso que se aplica un conjunto de prácticas para trabajar en equipo y de esta manera obtener el mejor resultado posible en un proyecto.

Se realizan entregas parciales hasta realizar la entrega del producto final. Estas entregas parciales las denominamos Sprints, al final de cada Srpint el equipo se reúne y decide cuales son los objetivos de cara al próximo Sprint.

Scrum es indicado para aquellos proyectos en los que  tenemos requisitos poco definidos, que están abiertos cambios, ya que lo indicado entonces es realizar objetivos a corto plazo y así  agilizar el desarrollo del mismo.

\subsection{Web scraping}

Técnica utilizada para simular  la navegación de un humano en Internet con el fin de extraer información del sitio web. Los datos pueden ser almacenados y analizados en una base de datos central o en otros lugares de almacenamiento.  Son diversas las técnicas que se pueden utilizar, en concreto hemos realizado a través de peticiones al protocolo HTTP. Hay algunas herramientas que nos permiten la extacción de datos como 80legs y Scrapinghub. Para el aprendizaje se ha decidido llevarlo a cabo completamente a mano.

\section{Herramientas}

\subsection{Trello}

Trello \url{https://trello.com} es una herramienta colaborativa que nos permite organizar nuestro  proyecto en tablones. En el tablón ponemos diferentes tarjetas y en cada una de ellas una lista.
Hemos decidido usarla principalmente por su sencillez y facilidad de uso. Otra razón que nos ha llevado a su uso es la aplicación para Android que nos permite estar  conectados en cualquier momento.

\subsection{GitHub}
Git Hub \url{https://github.com/} es una plataforma que nos permite el desarrollo colaborativo del software, alojando el repositorio de código en ella.
Su elección se debe al conocimiento de su funcionamiento, además de su buena adaptación a Ubuntu, sistema operativo en el que se ha trabajado y que mendiante sencillos comandos nos permite utilizarla fácilmente. El enlace a mi repositorio es el siguiente \url{https://github.com/nunobasurto/apuestas_deportivas}

\subsection{Oracle VM Virtual Box}
Virtual Box \url{https://www.virtualbox.org/} es un software libre con el cual hemos trabajado a lo largo de toda la carrera, cumple bien nuestras necesidades además de ser sencillo en su uso. Hay algunas alternativas como por ejemplo Hyper X que se pueden instalar en Windows, pero es necesario tener el sistema operativo en una versión Pro a la cual no tenemos acceso.

\subsection{Servidor web Apache}
Apache \url{http://httpd.apache.org/} es un servidor web Open Source y multiplataforma, utilizado para realizar servicio a paginas web estáticas o dinámicas.
Su elección de debe a que es el servidor web más conocido, además de tratarse de un software Open Source. Se adapta perfectamente a nuestras necesidades en el trabajo de fin de grado.


\subsection{Drupal}
Drupal \url{https://www.drupal.org/} es un gestor de contenido Open Source que posee una gran comunidad, es combinable con MySQL. Posee una gran cantidad de módulos sobre los que apoyarse.
Su sencillez de combinación con MySQL y el hecho de que se trata de un software Open Source, nos han hecho decantarnos por este sistema gestor de contenido. La gran cantidad de módulos también han contribuido a su uso.
La versión elegida ha sido Drupal 7 dado que esta tiene una madurez mucho mayor que Dupal 8, sobre todo para la elaboración de una pequeña web como la nuestra, la falta de tiempo para aventurarse y la comunidad detrás han sido claves.

\subsection{Sublime Text}

Sublime Text \url{https://www.sublimetext.com/} nos permiten editar código fuente de un programa, ayudando en la simplificación de la escritura y resaltando la sintaxis haciendo mas sencillo escribir el código es un editor de texto gratuito, que no libre que nos permite trabajar con una gran variedad de idiomas. En concreto nos permite trabajar con PHP, que ha sido el lenguaje en el que hemos desarrollado gran parte del proyecto.

\subsection{Zotero}

Zotero \url{https://www.zotero.org/} es una herramienta de gestión de información que  nos ayuda a gestionar las referencias bibliográficas. Obetenemos las referencias que deseamos utilizando la extensión para el navegador Chrome y lo exportamos en formato Bib\TeX a \LaTeX.

\subsection{Texmaker}

Texmaker \url{http://www.xm1math.net/texmaker} es una moderna plataforma que integra as diferentes herramientas que se ncesitan para desarrollar documentos con \LaTeX. Su uso viene motivada por tratarse de una de las mejores herramientes para \LaTeX y que tiene una licencia GPL.

\capitulo{5}{Aspectos relevantes del desarrollo del proyecto}


\section{Lenguaje para el desarrollo del proyecto}

Desde un primer momento se decidió el empleo de PHP como lenguaje central en la implementación del proyecto dado que permite el desarrollo web de contenido dinámico. El proyecto se ha llevado a cabo con Sublime Text, lo cual ha permitido una implementación del código más sencilla ya que iba mostrando posibles predicciones de texto.

PHP ha hecho posible aplicar alguna técnica de programación orientada a objetos, que se puede observar en el código del algoritmo de Backpropagation donde existen dos clases.
Se barajó la opción de implementar el código en Python y dejarlo embebido sobre PHP pero la idea inicial se descartó debido a que algunas funciones como las de acceso a la base de datos son únicas desde Drupal.

\subsection{PHP en Drupal}
La base de datos de MySQL estaba vinculada con Drupal, por lo que existe acceso desde Drupal a la base de datos. Para el acceso, el código PHP debía tener algunas sentencias exclusivas. La función db\_query tiene una sintaxis similar a las que se conocen ya que el contenido es una consulta MySQL sencilla. En cambio, otras funciones como db\_update, db\_select o db\_insert tienen una estructura completamente diferentes a las conocidas aunque son sencillas de utilizar.

Para poder desarrollar todo el código sobre Sublime Text y que no hubiese necesidad de copiarlo continuamente en los nodos de Drupal, fue necesario añadir un inclue\_once tanto en el nodo de Drupal como en el script.

\section{Datos de entrada a la red neuronal}
Es preciso darle un sentido a los datos que se han ido recopilando mediante los algoritmos de scraping, ya que el objetivo es darle a la red neuronal unos datos a partir de los cuales pueda ir aprendiendo y encontrar patrones. Son dos los algoritmos de scraping que fueron recopilando los datos que más tarde han sido empleados como entrada a la red Neuronal.

El primer algoritmo extrae todas las estadísticas de cada uno de los partidos de una jornada y los almacena en la base de datos. El segundo extrae la posición y las estadísticas de cada equipo en cada jornada, para almacenarlos posteriormente en la base de datos. A partir de estos datos recopilados se calculan las rachas de los equipos, donde se tiene en cuenta los puntos de cada equipo en las últimas jornadas acontecidas, los goles a favor y  los goles en contra. De esta forma se favorece al algoritmo para conocer la tendencia del equipo en las últimas jornadas. 

En el training, por cada una de las instancias se le pasan al algoritmo un total de 78 columnas, donde las 30 primeras contienen los estadísticas del partido, de la 30 a la 54 las rachas y las estadísticas del equipo local en la clasificación y de la 54 a la 78 lo correspondiente al equipo visitante. Como función objetivo se establece 0 si ganó el equipo local, 1 si ganó el equipo visitante y 0.5 si el resultado final fue empate.
Antes de la entrada de los datos al algoritmo se realiza una normalización de los mismos respecto a la misma columna del resto de instancias, siendo la normalización:
\begin{center}
\large{$$v' = \frac{v-min_{a}}{max_{a}-min_{a}}$$}
\end{center}


\section{Optimización del algoritmo de Backpropagation}
Cuando los datos se encuentran en la red neuronal es preciso optimizar el número de neuronas que se van a utilizar. Con un número de neuronas pequeño, la red neuronal puede ser incapaz de aprender todos los patrones existentes, teniendo problemas para devolver el ouput deseado. En cambio, si el número de neuronas es demasiado grande, la red neuronal se queda sin margen de adaptación a cambios.

Otro factor a tener en cuenta son las iteraciones que va a ejecutarse el training, ya que ante un pequeño número el algoritmo puede no aprender suficiente, y si el número es muy alto termina memorizando cada una. Esto supone que no se logren los resultados deseados en el test, por lo que se ha llevado a cabo un estudio para determinar cual es la cantidad más adecuada de neuronas y de epochs para el aprendizaje de la red neuronal.
De esta manera, se considera la mejor  configuración para la red neuronal aquella que menos  error ha dado.

\section{Interfaz sencilla para el usuario}

Con todos los datos necesarios es importante conocer cómo mostrarlos al usuario que accede al sitio web. Dado que el objetivo de esta página es su uso con apuestas deportivas, se ha hecho un algoritmo de scraping que recopila las cuotas de las casas de apuestas para cada partido.
A partir de los resultados que se han predicho, se muestran las cuotas de las cuatro casas de apuestas principales en España, ofreciendo al usuario la posibilidad de apostar en ellas.
En una de la páginas del sitio web se puede observar el balance general a lo largo de la temporada, donde es posible acceder a cada una de las jornadas y ver más detalladamente lo ocurrido a lo largo de la temporada.

\section{Automatización del funcionamiento}

Dado que la obtención de datos debe hacerse periódicamente y la ejecución manual por parte del usuario es costosa, se ha automatizado la ejecución de todos los algoritmos en base al momento en el que deban ejecutarse. Por ejemplo, solo es posible obtener los resultados de una jornada una vez finalizada esta, o es posible fijarse en las cuotas de las casas de apuestas, las cuales pueden ir variando a lo largo de la semana debido a que no son valores fijos.
Para cada jornada se establece una fecha\_antes y una fecha\_despues, los algoritmos de scraping de resultados se ejecutan una vez finalizada la jornada, es decir, si son posteriores a fecha\_despues, algo similar ocurre con el scraping de casas de apuestas solo que en este caso lo mejor es obtener los datos justo antes de la jornada para tener los datos más recientes. Para el algoritmo de backpropagation no es necesario el día exacto en el que ejecutarse, pero debe ejecutarse siempre una vez finalizada la jornada anterior, ya que son necesarias las rachas y los datos de la pasada jornada.

Todo esto lo logramos gracias al uso de demonios, también llamados servicios, que comparan la fecha actual a la fecha de la jornada, ya sea la anterior o posterior, y en caso de coincidir ejecutan el algoritmo deseado. El demonio se ejecuta todos los dias a la misma hora esto se debe a que no todas las jornadas comienzan el viernes y acaban el Lunes si no que algunas jornadas se disputan entre semana.
\capitulo{6}{Trabajos relacionados}

Se va a hablar otros trabajos que utilizan técnicas de Machine Learning en la predicción de apuestas deportivas, como vamos a ver no solo se llevan a cabo en fútbol.

\section{Neural Network Prediction of NFL Football Games}

Este trabajo de Joshua Kahn trata de la predicción de futuros partidos en la NFL.
El utiliza el lenguaje Perl y Matlab para la red neuronal. Los datos que utiliza lo extrae de la página oficial de la NFL con hojas de cálculo, en mi proyecto en cambio se están utilizando técnicas de web scraping y el almacenamiento en bases de datos.

J. Kahn utiliza técnicas de aprendizaje supervisado, más en concreto backpropagation, lo mismo que yo~\cite{neural_nfl}.

Aunque tiene un enfoque diferente a la hora de insertar datos y el deportes es diferente, el objetivo del proyecto y las técnicas utilizadas hacen que sea similar al mío.

\section{Sistema de predicción de resultados en eventos deportivos}

Trabajo de fin de grado llevado a cabo por Fernando Valera Guardiola, en la universidad Carlos III de Madrid. Ha utilizado WEKA para la creación de modelos de predicción, también ha utilizado un algoritmo genético para que la combinación de apuestas, maximice el beneficio minimizando el riesgo.

La recopilación de datos no fue capaz de realizarla de manera automática, por lo que tuvo que hacerla manualmente, algo muy tedioso ya que se invierte una gran cantidad de tiempo. Por contra el proyecto que he desarrollado recopila los datos mediante web scraping, además de ser automático gracias al demonio de Linux.

Sin duda es un enorme trabajo el que llevó a cabo que no solo predice los partidos de la liga, si no también de otras grandes ligas europeas y de otros deportes como el baloncesto donde tiene en cuenta la NBA. También destacar que el tiempo de realización del trabajo han sido 8 meses~\cite{apuestas_deportivas}.
\capitulo{7}{Conclusiones y Líneas de trabajo futuras}

Todo proyecto debe incluir las conclusiones que se derivan de su desarrollo. Éstas pueden ser de diferente índole, dependiendo de la tipología del proyecto, pero normalmente van a estar presentes un conjunto de conclusiones relacionadas con los resultados del proyecto y un conjunto de conclusiones técnicas. 
Además, resulta muy útil realizar un informe crítico indicando cómo se puede mejorar el proyecto, o cómo se puede continuar trabajando en la línea del proyecto realizado. 

Los objetivos que se tenían al comienzo del proyecto se han cumplido ya que hemos logrado una interfaz intuitiva, la automatización del algoritmo y el principal objetivo el pronóstico de resultados de fútbol.

\section{Conclusiones}

El proyecto que se ha llevado a cabo me parece muy ambicioso, el número de técnicas que abarca es bastante grande, gracias a ello he podido aprender bastante.
Me hubiese gustado disponer de más tiempo para realizar estudios con una mayor exactitud ello hubiese permitido mejores resultados en la red neuronal. Pese a ello me encuentro muy satisfecho del trabajo llevado a cabo, ya que teniendo en cuenta el tiempo disponible me parece que ha sido un buen proyecto.

Los conocimientos adquiridos durante la realización de la carrera universitaria, me han ayudado de gran manera a llevar a cabo el proyecto. Las principales asignaturas sobre las que me he apoyado son:

\begin{itemize}
\item \textbf{Computación neuronal y evolutiva:}Sin duda la asignatura con la que guarda más relación ya que todos los conceptos de redes neuronales, así como el algoritmo de backpropagation han sido extraídos de aquí.

\item \textbf{Sistemas operativos: } La navegación por Ubuntu mediante líneas de comandos, así como la creación de demonios son conocimientos vistos en esta asignatura.

\item \textbf{Bases de datos: }El funcionamiento de la base de datos y las diferentes consultas llevadas a cabo, no hubiesen sido posible realizarlas sin los conocimientos tratados en esta asignatura. Aunque vimos PostgreSQL y aquí se ha trabajado con MySQL las similitudes son grandes.

\item \textbf{Gestión de proyectos: } Sin duda otra de las asignaturas fundamentales en la planificación y gestión del proyecto, las metodologías ágiles han sido necesarias.

\end{itemize}

En cuanto a los conocimientos de la carrera que se han echado un poco de menos, una asignatura en la que aprender programación web sobre la cual apoyarnos para el desarrollo de los scripts en PHP o el web scraping.
\section{Líneas de trabajo futuras}
La escalabilidad que posee este proyecto es muy grande, sin duda puede continuarse en un entorno de apuestas deportivas incluyendo otros deportes, o es más en el propio fútbol cabe la posibilidad de pronosticar una mayor cantidad de tipos de apuestas como por ejemplo los goles a favor, los coners etc...

\subsection{Incremento de la información}
En lo primero que trabajaría serían en ofrecer al algoritmo una mayor cantidad de información con la que aprender, como los jugadores disponibles que tiene o los partidos de descanso que ha tenido el equipo si se ha disputado una competición entre semana. 

Sin duda hay más factores a valorar que no han podido tenerse en cuenta, ya sea por falta de tiempo o por falta de recursos para recopilar esos datos.

\subsection{Ampliación del entorno de actuación}
El trabajo es ampliable a otros depotes como el baloncesto o el balonmano , sin llevar a cabo una gran modificación del mismo. Sin duda hubiese sido atráctivo abarcar un mayor número de ligas europeas así como las competiciones europeas.

Fuera del fútbol, hubiese sido atractivo tratar con la NBA, sin duda la liga de baloncesto con mayor número de seguidores en el mundo, donde la gran cantidad de partidos que se disputan nos darían muchos datos con los que trabajar.

\subsection{Pronóstico en tiempo real}
Sin duda se trata de una de las tareas más complejas y a la vez más atractivas sobre las que trabajar en este proyecto. Que el usuario pueda acceder en tiempo real durante la disputa de partido a los pronósticos, supone un cambio en el proyecto pero seguiría los objetivos marcados por el mismo.


\bibliographystyle{plain}
\bibliography{bibliografia}

\end{document}